\section{機械学習に現れる線形代数計算}
\label{sec:linear-regression-normal-equation}

本節では，主資料\cite{shimada2020}の線形回帰の説明を基に，
最小二乗法から線形方程式が得られる計算を補う．

\subsection{線形回帰の表し方}
\label{subsec:linear-regression-formulation}

訓練データの個数を$m$とし，データ集合を
\begin{align*}
  \mathcal{D}
  \coloneq
  \left\{(\bm{x}_i,y_i)\mid i=1,2,\ldots,m\right\}
\end{align*}
とする．
ここで，$\bm{x}_i\in\mathbb{R}^{d}$は$d$個の特徴量を持つ入力ベクトル，
$y_i\in\mathbb{R}$は対応する目的変数である．
線形回帰では，回帰係数ベクトル
\begin{align*}
  \bm{w}
  \coloneq
  \begin{pmatrix}
    w_1 & w_2 & \cdots & w_d
  \end{pmatrix}^{\mathsf{T}}
\end{align*}
を用いて，$i$番目の予測値を
\begin{equation}
  \widehat{y}_i
  \coloneq
  \bm{x}_i^{\mathsf{T}}\bm{w}
  \label{eq:linear-regression-prediction}
\end{equation}
と表す．

切片を含める場合は，各入力ベクトルの先頭に定数$1$を追加し，
対応する回帰係数を$\bm{w}$へ含めればよい．
以下では，この処理が既に行われているものとして計算する．

訓練データをまとめた計画行列と目的変数ベクトルを，それぞれ
\begin{align}
  X
  &\coloneq
  \begin{pmatrix}
    \bm{x}_1^{\mathsf{T}}\\
    \bm{x}_2^{\mathsf{T}}\\
    \vdots\\
    \bm{x}_m^{\mathsf{T}}
  \end{pmatrix}
  \in\mathbb{R}^{m\times d},
  \label{eq:design-matrix}\\
  \bm{y}
  &\coloneq
  \begin{pmatrix}
    y_1 & y_2 & \cdots & y_m
  \end{pmatrix}^{\mathsf{T}}
  \in\mathbb{R}^{m}
  \label{eq:target-vector}
\end{align}
と定義する．
このとき，全訓練データに対する予測値ベクトルは$X\bm{w}$となる．

\subsection{最小二乗法}
\label{subsec:least-squares-method}

線形回帰では，予測値と目的変数の差を残差という．
ここでは，残差の二乗和を最小にする回帰係数を求める．
損失関数を
\begin{equation}
  \mathcal{L}(\bm{w})
  \coloneq
  \|X\bm{w}-\bm{y}\|_2^2
  \label{eq:least-squares-loss}
\end{equation}
と定義する．
式\eqref{eq:least-squares-loss}を展開すると
\begin{align*}
  \mathcal{L}(\bm{w})
  =
  \bm{w}^{\mathsf{T}}X^{\mathsf{T}}X\bm{w}
  -2\bm{w}^{\mathsf{T}}X^{\mathsf{T}}\bm{y}
  +\bm{y}^{\mathsf{T}}\bm{y}
\end{align*}
となる．
したがって，$\bm{w}$に関する勾配は
\begin{align*}
  \nabla_{\bm{w}}\mathcal{L}(\bm{w})
  =
  2X^{\mathsf{T}}X\bm{w}-2X^{\mathsf{T}}\bm{y}
\end{align*}
である．
損失関数を最小にする点ではこの勾配が零になるため，
\begin{equation}
  X^{\mathsf{T}}X\bm{w}
  =
  X^{\mathsf{T}}\bm{y}
  \label{eq:normal-equation-original}
\end{equation}
を得る．
式\eqref{eq:normal-equation-original}を\textbf{正規方程式}という．

\subsection{線形方程式への書き換え}
\label{subsec:reduction-to-linear-system}

行列$A$とベクトル$\bm{b}$を
\begin{align}
  A
  &\coloneq
  X^{\mathsf{T}}X,
  \label{eq:regression-matrix-a}\\
  \bm{b}
  &\coloneq
  X^{\mathsf{T}}\bm{y}
  \label{eq:regression-vector-b}
\end{align}
と定義すると，正規方程式は
\begin{equation}
  A\bm{w}=\bm{b}
  \label{eq:regression-linear-system}
\end{equation}
となる．
計画行列$X$の列が線形独立であれば，$A$は実対称正定値行列であり，逆行列を持つ．
この場合，最小二乗解は
\begin{align*}
  \bm{w}
  =
  A^{-1}\bm{b}
  =
  (X^{\mathsf{T}}X)^{-1}X^{\mathsf{T}}\bm{y}
\end{align*}
で与えられる．

計画行列$X$の列が線形独立でない場合には，$A$が逆行列を持たないことがある．
また，互いに強く関係する特徴量があると，小さな誤差が解へ大きく影響することがある．
前者の行列を特異行列，後者のような行列を悪条件行列という．
この問題を和らげる方法の一つがリッジ回帰である．
リッジ回帰では，係数$\bm{w}$が大きくなりすぎないように
$\lambda\|\bm{w}\|_2^2$を加える．この項を正則化項といい，
$\lambda>0$を正則化係数という．損失関数を
\begin{equation}
  \mathcal{L}_{\lambda}(\bm{w})
  \coloneq
  \|X\bm{w}-\bm{y}\|_2^2
  +\lambda\|\bm{w}\|_2^2
  \label{eq:ridge-loss}
\end{equation}
と定義すると，正規方程式は
\begin{align*}
  (X^{\mathsf{T}}X+\lambda I)\bm{w}
  =
  X^{\mathsf{T}}\bm{y}
\end{align*}
となる．
このとき
\begin{equation}
  A_{\lambda}
  \coloneq
  X^{\mathsf{T}}X+\lambda I
  \label{eq:regularized-regression-matrix}
\end{equation}
は正定値行列となるため，HHLアルゴリズムに必要な逆行列を持つ．
ただし，$\lambda$を大きくすると回帰係数はより小さくなる．
そのため，計算の安定性と予測精度の両方を考えて$\lambda$を決める必要がある．
